\chapter{INTRODUCTION AND LANGUAGE DESIGN}

\section{Background and Problem Statement}

Propositional logic is a foundational component of mathematical logic and computer science, dealing with propositions that are either \textbf{true} or \textbf{false} and the connectives used to combine them~\citep{ben2012mathematical}. This report presents a \textbf{Propositional Logic Evaluator} that accepts expressions composed of truth values (\texttt{t}, \texttt{f}) and logical connectives ($\wedge$, $\vee$), and produces two outputs: the \textbf{evaluated truth value} and an equivalent expression in \textbf{prefix (Polish) notation}. Parenthesised grouping overrides default precedence, and $\wedge$ has higher priority than $\vee$. For example, \texttt{t $\vee$ f $\wedge$ f} evaluates to \texttt{t} with prefix notation \texttt{($\vee$ t ($\wedge$ f f))}.

The implementation applies compiler front-end techniques~\citep{aho2006compilers}: lexical analysis, LALR(1) parsing, AST construction, and tree-based evaluation. The GUI is built with PySide6 (Qt6 for Python) and supports both Unicode ($\wedge$, $\vee$) and ASCII (\texttt{\&}, \texttt{|}) operator input. Negation ($\neg$), implication ($\rightarrow$), and biconditional ($\leftrightarrow$) are outside the scope of this assignment.

\section{Propositional Logic}

The language of propositional logic expressions is defined \emph{recursively} (as specified in the assignment):

\begin{enumerate}
    \item Any of the truth values \texttt{t} (true) and \texttt{f} (false) is a propositional logic expression.
    \item If $B$ and $B'$ are propositional logic expressions, then $B \wedge B'$ and $B \vee B'$ are propositional logic expressions.
\end{enumerate}

Additionally, parenthesised grouping is supported: if $B$ is an expression then $(B)$ is an expression, allowing the user to override default precedence.

\begin{table}[H]
\centering
\caption{Truth tables for conjunction ($\wedge$) and disjunction ($\vee$)}
\label{table:truth_tables}
\begin{tabular}{@{}cccc@{}}
\toprule
$p$ & $q$ & $p \wedge q$ & $p \vee q$ \\
\midrule
t & t & t & t \\
t & f & f & t \\
f & t & f & t \\
f & f & f & f \\
\bottomrule
\end{tabular}
\end{table}

\textbf{Conjunction} ($\wedge$): The result is \texttt{t} only when \emph{both} operands are \texttt{t}. \\
\textbf{Disjunction} ($\vee$): The result is \texttt{t} when \emph{at least one} operand is \texttt{t}.

\section{Alphabet and Token Specification}

The input language is defined over the alphabet shown in Table~\ref{table:tokens}. Each symbol or symbol class is associated with a token type produced by the lexical analyser.

\begin{table}[H]
\centering
\caption{Token specification for the propositional logic language}
\label{table:tokens}
\begin{tabular}{@{}llll@{}}
\toprule
Token Name & Regular Expression & Matches & Description \\
\midrule
\texttt{TRUE}   & \verb|t|       & \texttt{t}             & Boolean true literal \\
\texttt{FALSE}  & \verb|f|       & \texttt{f}             & Boolean false literal \\
\texttt{AND}    & \lstinline![∧&]!   & $\wedge$ or \texttt{\&} & Conjunction operator \\
\texttt{OR}     & \lstinline![∨|]!  & $\vee$ or \texttt{|}    & Disjunction operator \\
\texttt{LPAREN} & \verb|\(|     & \texttt{(}              & Left parenthesis \\
\texttt{RPAREN} & \verb|\)|     & \texttt{)}              & Right parenthesis \\
\bottomrule
\end{tabular}
\end{table}

Both Unicode symbols ($\wedge$, $\vee$) and ASCII alternatives (\texttt{\&}, \texttt{|}) are accepted. Whitespace (spaces and tabs) is silently ignored.

\section{Context-Free Grammar}

The syntax of the propositional logic language is defined by the following context-free grammar (CFG):

\begin{equation}
G = (V_N,\; V_T,\; P,\; S)
\end{equation}

where:
\begin{itemize}
    \item $V_N = \{\; \text{statement},\; \text{expr}\; \}$ --- non-terminal symbols
    \item $V_T = \{\; \texttt{TRUE},\; \texttt{FALSE},\; \texttt{AND},\; \texttt{OR},\; \texttt{LPAREN},\; \texttt{RPAREN}\; \}$ --- terminal symbols
    \item $S = \text{statement}$ --- start symbol
    \item $P$ --- production rules, defined below
\end{itemize}

The production rules are:

\begin{table}[H]
\centering
\caption{Production rules of the context-free grammar}
\label{table:grammar_rules}
\begin{tabular}{@{}cll@{}}
\toprule
Rule & Production & Description \\
\midrule
(1) & $\text{statement} \rightarrow \text{expr}$                        & Top-level entry point \\
(2) & $\text{expr} \rightarrow \text{expr}\;\texttt{OR}\;\text{expr}$   & Disjunction \\
(3) & $\text{expr} \rightarrow \text{expr}\;\texttt{AND}\;\text{expr}$  & Conjunction \\
(4) & $\text{expr} \rightarrow \texttt{LPAREN}\;\text{expr}\;\texttt{RPAREN}$ & Parenthesised grouping \\
(5) & $\text{expr} \rightarrow \texttt{TRUE}$                           & True literal \\
(6) & $\text{expr} \rightarrow \texttt{FALSE}$                          & False literal \\
\bottomrule
\end{tabular}
\end{table}

\textbf{Ambiguity.} Rules (2) and (3) make the grammar ambiguous because the expression \texttt{t $\vee$ f $\wedge$ f} admits two distinct parse trees: one where $\vee$ is the root (grouping as $\texttt{t} \vee (\texttt{f} \wedge \texttt{f})$) and one where $\wedge$ is the root (grouping as $(\texttt{t} \vee \texttt{f}) \wedge \texttt{f}$). This ambiguity is resolved by declaring operator \textbf{precedence} and \textbf{associativity} in the LALR(1) parser (Section~\ref{sec:precedence}).

\section{Operator Precedence and Associativity}
\label{sec:precedence}

To disambiguate the grammar, precedence levels and associativity are declared as shown in Table~\ref{table:precedence}.

\begin{table}[H]
\centering
\caption{Operator precedence and associativity}
\label{table:precedence}
\begin{tabular}{@{}ccc@{}}
\toprule
Precedence Level & Operator & Associativity \\
\midrule
1 (lower)  & $\vee$ (OR)  & Left \\
2 (higher) & $\wedge$ (AND) & Left \\
\bottomrule
\end{tabular}
\end{table}

\textbf{Consequence:} $\wedge$ binds tighter than $\vee$. In the expression \texttt{t $\vee$ f $\wedge$ f}:
\begin{equation}
\texttt{t} \vee \texttt{f} \wedge \texttt{f} \;\equiv\; \texttt{t} \vee (\texttt{f} \wedge \texttt{f}) \;\neq\; (\texttt{t} \vee \texttt{f}) \wedge \texttt{f}
\end{equation}

\textbf{Left associativity} means that chains of the same operator group from left to right:
\begin{equation}
\texttt{t} \vee \texttt{f} \vee \texttt{f} \;\equiv\; (\texttt{t} \vee \texttt{f}) \vee \texttt{f}
\end{equation}

These declarations are passed to the LALR(1) parser generator, which uses them to resolve shift/reduce conflicts in the generated parse table.

\section{Notation Systems}

The evaluator translates infix expressions into \textbf{prefix (Polish) notation}~\citep{lukasiewicz1957}, where the operator precedes its operands. Table~\ref{table:notations} compares the three standard notation systems.

\begin{table}[H]
\centering
\caption{Comparison of notation systems}
\label{table:notations}
\begin{tabular}{@{}llll@{}}
\toprule
Notation & Name & Example & Operator Position \\
\midrule
Infix   & Standard        & \texttt{t $\wedge$ f} & Between operands \\
Prefix  & Polish          & \texttt{($\wedge$ t f)} & Before operands \\
Postfix & Reverse Polish  & \texttt{t f $\wedge$}   & After operands \\
\bottomrule
\end{tabular}
\end{table}

The key advantage of prefix notation is that it eliminates the need for precedence rules --- the grouping structure is fully explicit through parenthesisation:
\begin{equation}
\text{Infix: } \texttt{t} \vee (\texttt{f} \wedge \texttt{f}) \;\;\Longrightarrow\;\; \text{Prefix: } (\vee\; \texttt{t}\; (\wedge\; \texttt{f}\; \texttt{f}))
\end{equation}
